% M3_NnaemekaNnachiEgwu.tex
% Milestone 3: Critical Evaluation and Transfer
% Real-Time Systems Seminar, HSHL, Prof. Dr. Stefan Henkler
% Author: Nnaemeka Nnachi-Egwu
%
% Structured notes format per seminar checklist.

\documentclass[conference,onecolumn]{IEEEtran}
\usepackage{amsmath, amssymb, booktabs, array, microtype, url}

\begin{document}

\title{Milestone 3: Critical Evaluation and Transfer\\
  {\large Priority Inheritance Protocol (Resource Access Protocols)}}

\author{\IEEEauthorblockN{Nnaemeka Nnachi-Egwu}
\IEEEauthorblockA{Electronic Engineering\\
Hochschule Hamm-Lippstadt\\
Lippstadt, Germany\\
nnaemeka.nnachi-egwu@stud.hshl.de}}

\maketitle

%---------------------------------------------------------------
\section{Concrete Strengths (with Buttazzo citations)}

\textbf{S1 --- Bounded priority inversion (Theorem~7.2, p.~219):}\\
Without any protocol, the inversion window is unbounded --- any number
of medium-priority tasks can prolong it indefinitely (Fig.~7.4).
PIP bounds blocking to at most $\alpha_i = \min(l_i, s_i)$ critical
sections per task. This is a \emph{formal guarantee}, not a heuristic.

\textbf{S2 --- Transparency (Buttazzo \S7.10, p.~247, Table~7.5):}\\
\texttt{pi\_wait}/\texttt{pi\_signal} are drop-in replacements for
standard semaphore primitives. Task \emph{application code} does not change.
No ceiling declarations required.
This is the only protocol (besides NPP) rated ``transparent.''
\emph{Practical consequence:} POSIX \texttt{PTHREAD\_PRIO\_INHERIT}
converts any existing mutex-protected code to PIP with one attribute
call. Mars Pathfinder fix: one line of configuration.

\textbf{S3 --- Low pessimism vs.\ HLP (Buttazzo p.~214):}\\
HLP blocks at activation time based on ceiling of \emph{entered} CS,
even if the resource-accessing branch is never taken at runtime.
PIP defers to actual \texttt{wait}: no blocking unless genuine
resource contention.
\emph{In practice:} PIP's observed blocking is usually below the
$B_i$ bound (the bound is worst-case, not average-case).

\textbf{S4 --- Transitive inheritance for nested CS (Lemma~7.2, p.~218):}\\
For nested critical sections, priority propagates transitively up the
chain, preventing intermediate holders from being preempted.
Without this, a 3-level nesting would still suffer priority inversion
at the intermediate level.

\textbf{S5 --- Direct feed into schedulability (Buttazzo \S7.9):}\\
$B_i$ values (Eqs.~7.9--7.11) plug directly into Liu-Layland (Eq.~7.19)
and RTA (Eq.~7.22) without modification to the scheduling framework.
PIP is the only protocol that provides this tight integration with
standard fixed-priority analysis.

%---------------------------------------------------------------
\section{Concrete Limitations and Counterexamples}

\subsection{L1: No Deadlock Prevention (Buttazzo \S7.6.5, Fig.~7.13, p.~226)}

\textbf{Minimal counterexample:}
Two tasks ($\tau_1$ high, $\tau_2$ low), two semaphores ($S_a$, $S_b$):
\begin{itemize}
  \item $\tau_1$: \texttt{wait}($S_a$); \texttt{wait}($S_b$); \ldots
  \item $\tau_2$: \texttt{wait}($S_b$); \texttt{wait}($S_a$); \ldots
\end{itemize}
Timeline: $t_1$: $\tau_2$ locks $S_b$; $t_2$: $\tau_1$ preempts; $t_3$:
$\tau_1$ locks $S_a$; $t_4$: $\tau_1$ tries $S_b$ $\to$ blocked, $\tau_2$
inherits $P_1$; $t_5$: $\tau_2$ tries $S_a$ $\to$ \textbf{deadlock}.

\textbf{Why PIP cannot prevent it:}
\texttt{pi\_wait} only boosts priority; it never denies a lock on a
free semaphore. The deadlock arises from a lock-\emph{order} cycle,
which PIP has no mechanism to detect.

\textbf{Buttazzo (p.~227):} ``The deadlock does not depend on the Protocol
but is caused by an erroneous use of semaphores.''

\textbf{Fix:} Total ordering on semaphore acquisition (by ID) eliminates
the cycle. Or replace PIP with PCP (Theorem~7.3: PCP prevents deadlock).

\textbf{Note:} PIP does not \emph{cause} deadlocks --- it is neutral.
The same deadlock would occur without any protocol.

\subsection{L2: No Chained Blocking Prevention (Buttazzo \S7.6.5, Fig.~7.12, p.~225)}

\textbf{Minimal counterexample:}
$\tau_1 \succ \tau_2 \succ \tau_3$; two semaphores $S_a$, $S_b$;
$\tau_2$ holds $S_b$; $\tau_3$ holds $S_a$.
$\tau_1$ arrives: blocked on $S_a$ (first blocking), then unblocked,
then blocked on $S_b$ (second blocking) $\to$ \textbf{two} blockings.

With $n$ distinct semaphores, each held by a different lower-priority task:
$\tau_1$ blocked up to $\alpha_1 = \min(n, n) = n$ times.
PCP (Theorem~7.4) limits this to exactly~1.

\subsection{L3: Kernel Implementation Complexity (Table~7.5, p.~248)}

PIP is the \emph{only} protocol rated ``Hard'' in Table~7.5.
Root cause: transitive chain traversal in \texttt{pi\_wait} must be
executed atomically in the kernel. Each step must:
(1) read current holder's blocked-on semaphore,
(2) boost that holder,
(3) repeat if that holder is also blocked.

This loop adds non-deterministic interrupt latency.
Production kernels (Linux \texttt{rtmutex.c}) bound chain depth
to prevent infinite loops --- which means formal guarantees may not
fully hold on the implementation.

Transparency (S2 above) is for \emph{application} programmers;
the \emph{kernel} implementation is decidedly non-trivial.

\textbf{Note on WCET sensitivity:}
$B_i$ formulas (Eqs.~7.9--7.11) require exact $\delta_{j,k}$ values.
Abella et~al.~[WCET] show no current WCET method is fully trustworthy
--- especially on modern hardware with caches, pipelines, and multicore
interference. Two distinct pessimism sources compound:
(1)~structural: $B_i^l$ sums CS that cannot simultaneously block
(Lemma~7.4 pessimism), (2)~input uncertainty: conservative WCET
estimates inflate $\delta_{j,k}$, making $B_i$ unnecessarily large.
\emph{This applies to all protocols with CS-duration-based blocking
formulas, not PIP specifically.}

%---------------------------------------------------------------
\section{Runtime, Memory, and Scalability}

\begin{itemize}
  \item Offline blocking computation: $\mathcal{O}(mn^2)$ per
        Fig.~7.11 algorithm. Feasible for $n, m \leq$ tens.
  \item Exact algorithm: $\mathcal{O}(2^n)$ --- impractical for
        $n > 15$.
  \item Runtime (per \texttt{pi\_wait}): chain traversal --- worst case
        proportional to nesting depth; bounded by implementation limits.
  \item Additional TCB fields: active priority, blocked-on semaphore ID,
        semaphore wait queues. $\mathcal{O}(n + m)$ extra state.
  \item Semaphore queues: $\mathcal{O}(n)$ total across all semaphores
        (each task in at most one queue at a time).
\end{itemize}

%---------------------------------------------------------------
\section{Transfer Analysis}

\subsection{Where PIP Is Suitable}

\begin{enumerate}
  \item \textbf{Single-core RTOS with bounded nesting:} $l_i$ and $s_i$
        are small; $\alpha_i$ is small; $B_i$ is manageable.
        Total semaphore ordering enforced $\to$ no deadlock risk.
  \item \textbf{POSIX legacy migration:} \texttt{PTHREAD\_PRIO\_INHERIT}
        is one attribute call; no code restructuring.
        Transparent, immediate improvement in worst-case timing.
  \item \textbf{Automotive ECU (classic AUTOSAR, single-core):}
        Bounded blocking time for ASIL certification;
        short CS WCETs for low-complexity ECU code;
        RM scheduling standard.
  \item \textbf{Any system where deadlock is provably absent by design:}
        If all semaphores acquired in fixed total order (common rule),
        PIP's only remaining weakness is chained blocking (Theorem~7.2
        bounds it), and the implementation overhead is acceptable.
\end{enumerate}

\subsection{Where PIP Is Unsuitable}

\begin{enumerate}
  \item \textbf{Multiprocessor/multicore systems:} Priority boost on
        a remote processor requires IPI; latency is non-deterministic.
        Gracioli et~al.~[GP]: IPI overhead has no equivalent in
        Buttazzo's $B_i$ formula $\to$ analysis invalid.
        Use MPCP, FMLP, or MSRP.
  \item \textbf{EDF-scheduled systems:} PIP undefined for dynamic priorities.
        Use SRP [Bak91] which natively supports EDF via preemption levels.
  \item \textbf{Safety-critical systems requiring certified deadlock freedom:}
        DO-178C, IEC~61508, ISO~26262 high-ASIL require
        provable deadlock absence. PIP alone is insufficient; use PCP
        or enforce and prove total semaphore ordering.
  \item \textbf{Deep nested CS with many tasks:} Transitive chain in
        \texttt{pi\_wait} adds atomic interrupt latency; $\alpha_i$ can
        be large; blocking bound grows.
  \item \textbf{Systems with unreliable WCET estimates:} If $\delta_{j,k}$
        cannot be bounded tightly (modern hardware, caches, multicore
        interference), the $B_i$ computation is unsound.
        $B_i$ under-estimates $\to$ false schedulability guarantee.
\end{enumerate}

%---------------------------------------------------------------
\section{Proposed Extensions}

\textbf{E1 --- POSIX validation:}
Implement the 3-task set with \texttt{pthread} + \texttt{SCHED\_FIFO},
run under \texttt{PTHREAD\_PRIO\_INHERIT} vs.\ plain mutex.
Measure actual $\tau_1$ blocking with \texttt{clock\_gettime}.
Hypothesis: observed blocking $\leq B_1 = 4$ in all runs,
and strictly less in most (bound is not tight).

\textbf{E2 --- UPPAAL deadlock witness:}
Add a second semaphore $S_a$ to the UPPAAL model with reverse
acquisition order in one task. Verify that \texttt{A[] not deadlock}
\emph{fails} and the counterexample trace matches Fig.~7.13,
making the deadlock limitation formally witnessed.

\textbf{E3 --- Tight blocking computation:}
Implement Rajkumar's exact $\mathcal{O}(2^n)$ algorithm in Python
for $n \leq 8$. Compare against Eq.~(7.11)'s result on the
4-task, 3-semaphore example (Buttazzo Table~7.1, p.~222)
to quantify the pessimism gap.

\textbf{E4 --- Multiprocessor overhead model:}
Using the overhead taxonomy from Gracioli et~al.~[GP] (Table~1:
$\Delta_\mathrm{cxs}$, $\Delta_\mathrm{scd}$, $\Delta_\mathrm{ipi}$),
derive a modified $B_i$ formula for a dual-core system with
Global-EDF + PIP semantics and check schedulability under
realistic IPI latency values from [GP].

\end{document}
